\documentclass[../../main.tex]{subfiles}

\begin{document}
\chapter{热力学基本规律}
本章作为整个热统笔记的开篇，将简要陈述热力学的基本概念和热力学第一第二定律，可以认为是对热学课程的复习。我并不打算以教科书的形式呈现内容，因此可能有罗列知识点的嫌疑，不过读者应该有过热学基础，不会感到内容的突兀。
\section{基本概念}
\subsection{热动平衡}
我们研究的热力学系统指的是包含大量微观粒子的宏观系统，在没有外界影响\footnote{倘若有外界影响，系统的各宏观性质也可能不随时间改变，这种状态称为稳态，比如稳恒电流、两侧温度固定的导热平板等}的条件下，当系统各部分的宏观性质不再随时间而改变，我们称系统处于热力学平衡态，简称平衡态，这是热统课程主要研究的系统\footnote{在本笔记中，提及系统大部分都是平衡态系统，请读者自行分辨}。因为系统的宏观性质不再改变，我们就可以测量系统的一些物理量，换而言之，我们可以用这些物理量描述一个平衡态，这些物理量即所谓的宏观量。确定了系统的宏观量，也就确定了系统的宏观状态，因此我们也称系统宏观量为态函数，即平衡态的函数。当然，因为系统处于平衡态这个强有力的约束，系统的独立宏观量往往很少，只需少量宏观量就可以描述系统的所有宏观量，这将在下一章深入学习。系统宏观量可以分为四类：几何变量、力学变量、电磁变量、化学变量，这在热学中已经学习过不再赘述。要注意的是，这些宏观量是描述系统整体的，也就是只对平衡态才有意义，因此对于一个实际的热力学过程，谈及过程中系统的状态是毫无意义的。正因如此，我们提出准静态过程的概念，这是假想出来的过程，这种过程进行的相当慢以至于系统偏离平衡态后又很快回归平衡态，因此在这个过程中，可以认为系统始终处于平衡态，可以用态函数描述，这在热学中重点论述过，此处不再赘述。



\end{document}
